\chapter{IMPLEMENTATION}

\section{Development Environment}
The project is developed as a monorepo using pnpm workspaces and Turborepo.
Development is on Windows 11 with WSL2 and Docker Desktop. VS Code is the
primary IDE. Git with Conventional Commits (enforced via Husky + Commitlint)
is used for version control.

\section{Tools and Technologies}
Table~\ref{tab:stack} summarizes the technology stack.

\begin{table}[htbp]
  \centering
  \caption{Technology Stack}
  \label{tab:stack}
  \begin{tabular}{@{}ll@{}}
    \toprule
    \textbf{Layer} & \textbf{Technology} \\
    \midrule
    Public Frontend       & Next.js 13, Tailwind CSS 3.4, Lucide \\
    Authenticated Frontends & React 18, Vite 5, Tailwind CSS 3.4 \\
    Backend Runtime       & Node.js 18 \\
    Backend Framework     & Express 4.18 \\
    AI Framework          & FastAPI 0.104, Python 3.10 \\
    AI Libraries          & Pandas, NumPy, scikit-learn \\
    Database              & PostgreSQL 14 \\
    ORM                   & Sequelize 6 \\
    Auth                  & JWT + bcrypt \\
    Containerization      & Docker, Docker Compose \\
    \bottomrule
  \end{tabular}
\end{table}

\section{Docker Setup}
All 13 containers run via \texttt{docker-compose.yml} on a custom bridge
network.

\section{Database Implementation}
Sequelize models are defined in catalog-service, auth-service, inventory-service,
and location-service. Transactions are used for multi-table operations.

\begin{lstlisting}[language=JavaScript, caption={Transaction example for laptop creation}, label={lst:transaction}]
const transaction = await sequelize.transaction();
try {
  const laptop = await Laptop.create({ brand, model, ... }, { transaction });
  await Specification.create({ laptop_id: laptop.id, ... }, { transaction });
  await transaction.commit();
} catch (error) {
  await transaction.rollback();
}
\end{lstlisting}

\section{Data Acquisition}
A Python scraper targets priceoye.pk, collecting 120 real laptops with PKR
prices.

\section{AI Recommendation Engine}
The engine uses weighted cosine similarity with dynamic feature weights.

\begin{lstlisting}[language=Python, caption={Weighted cosine similarity}, label={lst:cosine}]
def weighted_cosine_similarity(user_vector, laptop_vectors, weights):
    weighted_user = user_vector * weights
    weighted_laptops = laptop_vectors * weights
    dot = np.dot(weighted_laptops, weighted_user)
    norm_user = np.linalg.norm(weighted_user)
    norm_laptops = np.linalg.norm(weighted_laptops, axis=1)
    denom = norm_user * norm_laptops
    denom[denom == 0] = 1e-10
    return dot / denom
\end{lstlisting}

\section{API Gateway}
The API Gateway is the single public entry point with Helmet, CORS, rate
limiting, JWT verification, and admin IP restriction.

\section{Frontend Implementation}
Figure~\ref{fig:home-hero} shows the redesigned home page hero section.

\begin{figure}[htbp]
  \centering
  \includegraphics[width=0.9\textwidth]{figure-4-01-home-hero.png}
  \caption{Home Page -- Hero Section}
  \label{fig:home-hero}
\end{figure}

Figure~\ref{fig:auth-login} shows the redesigned authentication login page.

\begin{figure}[htbp]
  \centering
  \includegraphics[width=0.7\textwidth]{figure-4-06-auth-login.png}
  \caption{Authentication App -- Login Screen}
  \label{fig:auth-login}
\end{figure}

Figure~\ref{fig:customer-results} shows recommendation results.

\begin{figure}[htbp]
  \centering
  \includegraphics[width=0.9\textwidth]{figure-4-13-customer-results.png}
  \caption{Customer App -- Recommendation Results with WhatsApp Contact}
  \label{fig:customer-results}
\end{figure}

Figure~\ref{fig:vendor-inv} shows the vendor inventory dashboard.

\begin{figure}[htbp]
  \centering
  \includegraphics[width=0.9\textwidth]{figure-4-14-vendor-inventory.png}
  \caption{Vendor App -- Inventory Tab}
  \label{fig:vendor-inv}
\end{figure}

Figure~\ref{fig:admin} shows the admin dashboard.

\begin{figure}[htbp]
  \centering
  \includegraphics[width=0.9\textwidth]{figure-4-18-admin-dashboard.png}
  \caption{Admin Panel -- Dashboard with Stat Cards}
  \label{fig:admin}
\end{figure}

\section{Security Hardening}
Table~\ref{tab:security} lists the security layers.

\begin{table}[htbp]
  \centering
  \caption{Security Layers}
  \label{tab:security}
  \begin{tabular}{@{}ll@{}}
    \toprule
    \textbf{Layer} & \textbf{Implementation} \\
    \midrule
    HTTP headers     & Helmet on all Node services \\
    Authentication   & JWT (7-day expiry) \\
    Password storage & bcrypt (work factor 10) \\
    Rate limiting    & 100/min global, 20/15min auth \\
    Input validation & express-validator \\
    Sanitization     & Strips \$ and . keys \\
    SQL safety       & Sequelize parameterized queries \\
    Admin protection & IP allowlist + role check \\
    \bottomrule
  \end{tabular}
\end{table}

\section{CI/CD Pipeline}
GitHub Actions workflow builds all Docker images on every push to \texttt{main}.
